\section{A rational compiler for adaptive observation processes}
\label{sec:adaptivecompiler59}
The common-row construction also yields a causal approximation, but its
proof requires control of the full observation law. Correct terminal
means alone are not enough. We give that additional argument here.

The physical state is $x\in B_3$, initialized at a supplied rational
$x_0\in S^2$. A finite rational rotation list $\{A_a\}\subset SO(3)$
acts by $x\mapsto A_ax$ and emits a null symbol. A finite probe list has
rational affine probability laws
\begin{equation}\label{eq:affineprobes59}
 q_{j,y}(x)=b_{j,y}+\ip{\ell_{j,y}}x,\qquad
 \sum_y b_{j,y}=1,\quad \sum_y\ell_{j,y}=0,
 \quad b_{j,y}-\norm{\ell_{j,y}}\ge\eta>0.
\end{equation}
A probe emits a fresh draw and leaves $x$ unchanged. Controls may depend
on the entire observed transcript. A simulator has finite retained
labels and a nonnegative joint emission/transition matrix for each
control and output; these matrices may depend on horizon and epoch.
Let $\mathsf{PW}_{N,\varepsilon}$ denote its minimum peak label count
for one total-variation budget $\varepsilon$ on the entire transcript,
uniformly over all adaptive randomized policies. This paragraph defines
a causal model independently of the terminal matrix-output model.

Put $A=\max_j\sum_y\norm{\ell_{j,y}}^2$ and
$C_0=\max(3,\max_j(M_j-1))$. The lower probability bound in
\eqref{eq:affineprobes59} is used only for the relative-entropy upper
estimate, not for the finite-testing converse below.

\begin{theorem}[Adaptive rational and fair-bit realization]
\label{thm:adaptivecompiler59}
For rational $0<\varepsilon<1$ and integer $N\ge0$, choose
\begin{equation}\label{eq:adaptiveparameters59}
 m=\max\left(2,\left\lceil\sqrt{
                    \frac{4AN^2}{\eta\varepsilon^2}}\right\rceil\right),
 \qquad r=1-m^{-2},\qquad
 b=\max\left(0,\left\lceil\log_2
                 \frac{2C_0\max(1,N)}{\varepsilon}\right\rceil\right).
\end{equation}
There is an error-$\varepsilon$ causal realization on at most
$13m^2+1$ labels. Before probability rounding every rotation row is
rational, has at most four nonzero entries, and has barycentric mean
$rA_av$. Probe emissions leave the label unchanged. After rounding,
all emission and transition probabilities are dyadic, with at most $b$
fair bits per control and $Nb$ bits in total. The tables are stationary
within the chosen horizon.

An exact rational algorithm constructs these tables in time polynomial
in the numerical parameters $N,\varepsilon^{-1},A/\eta$, the explicit
input size, and the input bit lengths. In particular, for fixed probe
and rotation lists,
\begin{equation}\label{eq:adaptivewidth59}
                 \mathsf{PW}_{N,\varepsilon}
                    =O(1+N^2/\varepsilon^2).
\end{equation}
The bit count now includes the actual finite probe outputs as well as
hidden transitions; it does not describe quantum-state preparation.
\end{theorem}
\begin{proof}
Take the rational net $V_m$ from Lemma~\ref{lem:rationalnet58}, including
zero, and add $x_0$ if needed. This adds at most one label and cannot
shrink the convex hull. Initialize deterministically at $x_0$. For every
rotation and label $v$, the same exact triple search supplies a row with
mean $rA_av$. At probe $j$, emit $y$ with probability $q_{j,y}(v)$ and
retain $v$. All points lie in $B_3$.

Fix an arbitrary policy. On the same finite trajectory space of controls,
outputs and hidden labels, introduce an auxiliary law $P$: hidden labels
follow the just constructed rotation kernels, but every probe emits
$q_{j,y}(X)$ from the true physical state $X$, independently of the
hidden label conditional on the public history. Its public marginal is
exactly the target experiment. Let $Q$ be the constructed simulator law.
The policy kernels, initial label and hidden rotation kernels are the
same under $P,Q$; only probe emissions differ.

Under $P$, if the current physical and hidden points are $X,V$, both
$X$ and every nonzero hidden vertex are unit vectors. At a rotation
step, for $V\ne0$,
\begin{align*}
 \E[\norm{V'-A_aX}^2\mid X,V,a]
 &\le 2-2r\ip V X\\
 &=r\norm{V-X}^2+2(1-r).
\end{align*}
For $V=0$ choose its row to remain at zero; the same upper bound holds
because the left side is one and the right side is $2-r\ge1$.
The added seed is a unit vector, so these two cases exhaust the state
set. This is why a general interior seed has not been silently added
to the theorem. A probe changes neither point. Iteration, with the
control selected adaptively, gives
\begin{equation}\label{eq:processmoment59}
                  \E_P\norm{V_t-X_t}^2\le 2t/m^2.
\end{equation}
The inequality is unconditional; no posterior error bound after an
unlikely observation is assumed.

For probability vectors $p,q$ with $q_y\ge\eta$, the inequality
$\log z\le z-1$ gives
\[
 D(p\Vert q)\le\sum_y(p_y-q_y)^2/q_y.
\]
Consequently the conditional divergence at a probe is at most
$A\norm{X-V}^2/\eta$. The finite-trajectory relative-entropy chain rule
cancels the common policy and rotation kernels. Applying
\eqref{eq:processmoment59} at the probe times gives
\begin{equation}\label{eq:processKL59}
 D(P\Vert Q)\le \frac{2AN^2}{\eta m^2},\qquad
 \TV(P_{\rm pub},Q_{\rm pub})
       \le\sqrt{D(P\Vert Q)/2}
       \le \frac Nm\sqrt{A/\eta}\le\varepsilon/2.
\end{equation}
Here marginalization can only decrease divergence. These elementary
entropy facts apply to the joint classical trajectory, not to a quantum
state. One may verify the total-variation inequality by partitioning
into the set where $P\ge Q$ and its complement: the log-sum inequality
reduces divergence to a Bernoulli divergence, whose second derivative
in its first parameter is at least four. This yields
$D(P\Vert Q)\ge2\TV(P,Q)^2$. The chain rule follows by factoring each
trajectory probability into its conditional kernels. All probe
probabilities in $Q$ are positive, so absolute continuity causes no
exception. If $A=0$, every probe is constant and the divergence is zero.

Round each rotation row and each probe row by rounding all but its last
positive mass down to denominator $2^b$ and assigning the remainder to
the last. A row supported on $d$ points changes in total variation by
at most $(d-1)2^{-b}$. Thus every one-step joint kernel changes by at most
$C_0 2^{-b}$. Coupling the two label--transcript processes until their
first discrepancy proves, for every adaptive policy, an additional
path error at most $NC_0 2^{-b}\le\varepsilon/2$. Combining this with
\eqref{eq:processKL59} proves the asserted error. Each rounded row can
be sampled by one uniform $b$-bit integer and cumulative integer masses;
zero or deterministic rows may use fewer bits.

The net and an added rational unit seed are explicit. For every
command--vertex pair the exact enumeration in
Lemma~\ref{lem:rationalnet58} examines at most $\binom{k-1}{3}$ triples,
where $k\le13m^2+1$. The total is $O(|\mathcal A|k^4)$ constant-size
rational linear-algebra operations, in addition to the probe tables.
All determinants, comparisons and row identities involve rational
numbers with polynomially bounded bit lengths in the supplied data and
$\log m$; dyadic rounding adds $b$ bits. The computation of $m$ and $b$
can use integer comparisons, not floating-point decisions. This proves
the stated uniform, numerical-parameter complexity. It is a construction,
not a minimum-width algorithm, and its exhaustive fallback need not be
practical for large horizons.
\end{proof}

\subsection{An effective process converse for the rational free experiment}
Keep the rational rotations and seed of
\eqref{eq:rationaldata58}--\eqref{eq:rationalalphabet58}, now exposing the
three nondisturbing binary probes
\begin{equation}\label{eq:weakprobes59}
       q_{j,+}(x)=(1+\lambda x_j)/2,\qquad
       q_{j,-}(x)=(1-\lambda x_j)/2,
       \quad j=1,2,3,
\end{equation}
where $0<\lambda<1$ is rational. This defines a fixed rational classical
experiment. The probing law is not an implementation of nondisturbing
quantum measurement. Let $\mathsf{PW}$ refer to the causal model above,
not to terminal density-matrix width $W$.

\begin{theorem}[Effective repeated-observation lower bound]
\label{thm:processlower59}
Let $0\le\varepsilon<1$ and $0<\delta<1-\varepsilon$. For an integer
$t\ge0$ put
\begin{equation}\label{eq:testlength59}
 L_t=\max\left(1,\left\lceil18\lambda^{-2}625^t
                                      \log(6/\delta)\right\rceil\right).
\end{equation}
At every horizon $N\ge t+3L_t$, any error-$\varepsilon$ causal machine
has
\begin{equation}\label{eq:processcutlower59}
 |S_t|\ge(1-\varepsilon-\delta)(2\cdot3^t-1).
\end{equation}
If $\varepsilon+\delta<1/2$, it has the stronger lower bound
$|S_t|\ge2\cdot3^t-1$. More generally, the same bounds hold at every
cut $u$ with $t\le u\le N-3L_t$ after padding the prefixes by probes.
In particular, for each fixed $0<\varepsilon<1$, there are explicit
positive constants $c_{\varepsilon,\lambda},C_{\varepsilon,\lambda}$
such that, for all sufficiently large $N$,
\begin{equation}\label{eq:processgrowth59}
 c_{\varepsilon,\lambda}N^{\log3/\log625}
 \le\mathsf{PW}_{N,\varepsilon}
 \le C_{\varepsilon,\lambda}(N+1)^2.
\end{equation}
The converse uses no spectral-gap theorem. Its exponent is not asserted
to be sharp.
\end{theorem}
\begin{proof}
The rational freeness and stabilizer argument in
Theorem~\ref{thm:rationalseparation58} gives $2\cdot3^t-1$ different
reachable pure Bloch vectors at time $t$. Their coordinates lie in
$25^{-t}\Z^3$. Hence each distinct pair differs in some coordinate by
at least $25^{-t}$, and the associated plus-probabilities in
\eqref{eq:weakprobes59} differ by at least $\lambda25^{-t}/2$.
Use one common continuation: repeat each of the three probes $L_t$
times. Hoeffding's elementary Bernoulli bound and a union bound give
\[
 \mathbb P_x\left\{\max_j|\widehat q_j-q_{j,+}(x)|
                    \ge\lambda25^{-t}/6\right\}
 \le6\exp(-L_t\lambda^2 625^{-t}/18)\le\delta.
\]
The empirical neighborhoods for the different targets are disjoint.
At a cut label $z$, let $Q_z$ be the law of this same continuation.
Marginalizing prefix outputs and using total-variation correctness
shows that the continuation law from each prefix gives its own
neighborhood probability at least $1-\varepsilon-\delta$.
Writing this law as $\sum_z\alpha_x(z)Q_z$ and summing over the disjoint
neighborhoods yields \eqref{eq:processcutlower59}. If the probability
exceeds $1/2$, each neighborhood needs a different label assigning it
probability greater than $1/2$, giving the stronger conclusion.
The proof does not condition on the success of a previous observation.
Padding by arbitrary fixed probes leaves the target states unchanged
and proves the assertion at later cuts. Unused final controls may also
be probes. A terminal cut without enough remaining observations is not
claimed to satisfy this lower bound.

For the power lower bound take $\delta=(1-\varepsilon)/2$ and
$B=54\lambda^{-2}\log(6/\delta)+4$. Then
$t+3L_t\le B625^t$. For $N\ge B$, choose
$t=\lfloor\log_{625}(N/B)\rfloor$. Since
$3^t\ge(N/B)^{\log3/\log625}/3$ and $2\cdot3^t-1\ge3^t$,
\eqref{eq:processcutlower59} gives, for example,
$c_{\varepsilon,\lambda}=(1-\varepsilon)B^{-\log3/\log625}/6$.
The upper bound follows from Theorem~\ref{thm:adaptivecompiler59} with
$\eta=(1-\lambda)/2$ and $A=\lambda^2/2$. For irrational requested
$\varepsilon$, construct at any rational tolerance between
$\varepsilon/2$ and $\varepsilon$. All constants can thus be bounded
from the displayed data without computing a group spectral gap.
\end{proof}

The tests in this theorem consume horizon. The full-state terminal
exponential profile and the causal lower bound therefore express
different requirements. Conversely, the causal upper construction
controls one joint transcript law under adaptive control, not just its
one-time marginals. These process statements do not by themselves sharpen the terminal
spectral-gap converse. Its logarithmic loss is removed separately by
Theorem~\ref{thm:sharpoccupation61}.
